Chương mở đầu này cung cấp bối cảnh lý do chọn đề tài, 
xác định mục tiêu, đối tượng và phạm vi nghiên cứu, đồng thời 
trình bày cấu trúc tổng thể của đồ án. 
Đây là nền tảng để hiểu rõ ý nghĩa và định hướng của đồ án.
\section{Động cơ nghiên cứu}

Trong những năm gần đây, kiến trúc microservices và nền tảng container hóa đã trở thành xu hướng 
chủ đạo trong việc triển khai các hệ thống phân tán hiện đại. Kubernetes hiện là nền tảng 
điều phối phổ biến nhất, được sử dụng rộng rãi để quản lý container, tự động triển khai, mở rộng 
và vận hành hệ thống dịch vụ quy mô lớn. Một trong những đặc điểm quan trọng của Kubernetes là 
khả năng \textit{tự phục hồi} (self-healing), cho phép hệ thống tự động khởi động lại Pod lỗi, thay thế Pod bị crash hoặc điều phối workload sang Node khác khi xảy ra sự cố.

Tuy nhiên, cơ chế tự phục hồi mặc định của Kubernetes chủ yếu hoạt động dựa trên các rule cố định và các controller truyền thống. Trong các môi trường thực tế có tải động, số lượng dịch vụ lớn và lỗi xảy ra liên tục, các cơ chế này vẫn tồn tại nhiều hạn chế như:
\begin{itemize}
    \item Thời gian phục hồi chưa tối ưu.
    \item Không thích nghi tốt với nhiều loại lỗi đồng thời.
    \item Thiếu khả năng học hỏi từ dữ liệu lịch sử.
    \item Không tối ưu tài nguyên trong quá trình phục hồi.
\end{itemize}

Trong khi đó, Học tăng cường (Reinforcement Learning) là một nhánh của Học máy (Machine Learning) 
cho phép tác nhân (agent) học cách ra quyết định tối ưu thông qua quá trình tương tác với môi trường. 
Học tăng cường đã được ứng dụng trong nhiều bài toán tối ưu động như robot, mạng máy tính, quản lý 
tài nguyên đám mây và tối ưu lập lịch. Khả năng học chính sách tối ưu dựa trên trạng thái hệ thống 
giúp Học tăng cường trở thành hướng tiếp cận tiềm năng cho bài toán tự phục hồi trong Kubernetes.

Từ đó, đề tài “Ứng dụng Học tăng cường cho cơ chế tự phục hồi trong cụm Kubernetes” được thực hiện 
nhằm nghiên cứu khả năng áp dụng Học tăng cường để hỗ trợ ra quyết định phục hồi hệ thống hiệu quả 
hơn so với cơ chế mặc định.

\section{Mục tiêu đề tài}

Đề tài hướng đến các mục tiêu chính sau:
\begin{itemize}
    \item Nghiên cứu và phân tích các cơ chế tự phục hồi hiện tại của Kubernetes, đặc biệt trong 
    kịch bản Pod failure.
    \item Xây dựng môi trường ứng dụng microservices trên cụm Kubernetes. 
    \item Mô phỏng lượng người dùng và kịch bản lỗi thực tế.
    \item Thu thập dữ liệu trạng thái hệ thống thông qua các công cụ giám sát, phục vụ huấn luyện mô hình Học tăng cường.
    \item Thiết kế mô hình Học tăng cường hỗ trợ lựa chọn hành động phục hồi phù hợp.
    \item Đánh giá hiệu quả của mô hình sau huấn luyện dựa trên các tiêu chí như thời gian phục hồi, độ ổn định hệ thống 
    và khả năng thích nghi với lỗi.
\end{itemize}

\section{Phạm vi nghiên cứu}

Phạm vi nghiên cứu của đề tài tập trung vào bài toán tự phục hồi (self-healing)
trong môi trường Kubernetes thông qua việc ứng dụng học tăng cường nhằm hỗ trợ
quá trình ra quyết định phục hồi hệ thống khi xảy ra sự cố.

Nghiên cứu hướng đến việc xây dựng cơ chế phục hồi thông minh có khả năng
học từ dữ liệu vận hành của hệ thống để tối ưu hóa hiệu quả cứu hộ
trong môi trường microservices.

Đề tài được giới hạn trong kịch bản lỗi Pod CrashLoopBackOff. Đây là một
trong những sự cố phổ biến nhất, ảnh hưởng trực tiếp đến tính sẵn sàng
của hệ thống và thường gặp trong môi trường Kubernetes thực tế.

Môi trường thực nghiệm được xây dựng dưới dạng một cụm Kubernetes triển khai
ứng dụng microservices nhằm mô phỏng quá trình vận hành của hệ thống phân tán.

Dữ liệu được thu thập từ Kubernetes và các công cụ giám sát để xây dựng
tập dữ liệu huấn luyện cho bài toán Học tăng cường ngoại tuyến (Offline RL).

Tập dữ liệu bao gồm các thành phần cơ bản của mô hình như trạng thái hệ thống (state),
hành động phục hồi (action), phần thưởng (reward) và trạng thái kế tiếp (next state),
phản ánh rõ mối quan hệ giữa vận hành cụm và hiệu quả phục hồi.

Trong phạm vi nghiên cứu, nhóm dùng phương pháp huấn luyện dựa trên mô phỏng,
sau đó mới triển khai và đánh giá trên hệ thống thực nghiệm (Lab/Staging environment). 
Cách tiếp cận này giúp giảm nguy cơ gây mất ổn định và đảm bảo an toàn tuyệt đối
trước khi mô hình được xem xét cho các hệ thống thực tế.

Các hành động phục hồi được xem xét bao gồm: scale Pod, restart Pod,
reschedule workload, cordon hoặc drain Node, cũng như rollback dịch vụ khi xảy ra lỗi.

Hiệu quả của mô hình được đánh giá dựa trên thời gian phục hồi dịch vụ,
mức độ duy trì tính sẵn sàng, độ ổn định khi lỗi xảy ra liên tục
và khả năng thích nghi với nhiều trạng thái lỗi khác nhau.

Bên cạnh đó, đề tài không tập trung nghiên cứu các lĩnh vực ngoài bài toán tự phục hồi,
bao gồm: bảo mật và phát hiện tấn công, tối ưu autoscaling hoặc scheduling tổng quát,
điều phối đa cụm, cũng như triển khai trên môi trường production quy mô lớn.

Việc giới hạn này nhằm tập trung đánh giá sâu hiệu quả của Học tăng cường
đối với bài toán tự phục hồi trong Kubernetes một cách rõ ràng và nhất quán.

\section{Cấu trúc báo cáo}

Nội dung báo cáo được tổ chức thành bảy chương chính như sau:

\begin{itemize}
    \renewcommand{\labelitemi}{-}
    \item Chương 1: \textbf{Giới thiệu} -- Trình bày động cơ nghiên cứu, mục tiêu, phạm vi thực hiện và cấu trúc tổng thể của báo cáo.
    \item Chương 2: \textbf{Tổng quan lý thuyết, nghiên cứu liên quan} -- Tóm tắt các khái niệm nền tảng về Kubernetes, cơ chế self-healing, học tăng cường và các công trình liên quan đến tự phục hồi trong hệ thống cloud-native.
    \item Chương 3: \textbf{Mô hình bài toán} -- Mô tả bài toán self-healing dưới góc nhìn ra quyết định tuần tự, xây dựng mô hình MDP, không gian trạng thái, không gian hành động và hàm phần thưởng.
    \item Chương 4: \textbf{Phương pháp đề xuất} -- Trình bày kiến trúc phương pháp tiếp cận, quy trình thu thập dữ liệu, huấn luyện tác nhân RL và cách tích hợp mô hình vào bài toán tự phục hồi.
    \item Chương 5: \textbf{Cài đặt và thực nghiệm} -- Mô tả môi trường triển khai trên AWS, cụm k0s, hệ thống mini-ecommerce, công cụ giám sát, kịch bản tạo lỗi và cách triển khai tác nhân RL trên Kubernetes.
    \item Chương 6: \textbf{Kết quả và phân tích} -- Trình bày kết quả huấn luyện, kết quả đánh giá các chính sách, so sánh giữa các mô hình và phân tích tác động của cơ chế self-healing đến các chỉ số hệ thống.
    \item Chương 7: \textbf{Kết luận và hướng phát triển} -- Tóm tắt các kết quả đạt được, nêu ra hạn chế còn tồn tại và đề xuất các hướng mở rộng trong tương lai.
\end{itemize}

